\section{4. DISCUSSION}\label{discussion}

\subsection{4.1. Interpretation against research questions and hypotheses}\label{interpretation-against-research-questions-and-hypotheses}

\subsubsection{4.1.1. RQ1, Can a locally trained XGBoost model outperform persistence and CAMS?}\label{rq1-can-a-locally-trained-xgboost-model-outperform-persistence-and-cams}

The answer is partially affirmative. At the 6-hour horizon, XGBoost achieves RMSE = 23.44 μg/m³ and R² = 0.545, outperforming persistence by 22.2\% in RMSE and CAMS by 41.1\% in RMSE. Hypothesis H1 is therefore confirmed for the 6-hour horizon.

At the 12 and 24-hour horizons the picture is more nuanced. XGBoost outperforms both persistence and CAMS (which has negative R² at all horizons), confirming that locally trained models are more informative than either naive extrapolation or the global CTM at longer lead times. However, Ridge regression outperforms XGBoost at both longer horizons: by 9.8\% in RMSE at +12 h (Ridge: 25.85 μg/m³ vs.~XGBoost: 28.67 μg/m³) and by 13.5\% at +24 h (Ridge: 27.24 μg/m³ vs.~XGBoost: 31.48 μg/m³). Hypothesis H1 is therefore partially falsified for the 12 and 24-hour cases, where the best locally trained model is Ridge regression rather than XGBoost.

The margin of victory of locally trained models over CAMS is consistent across all three horizons and both model types. At +6 h, XGBoost RMSE is 23.44 vs. CAMS RMSE of approximately 39.7 μg/m³ (implied by R² = −0.314 on the same test set). At +12 h, Ridge RMSE is 25.85 vs.~CAMS approximately 43.1 μg/m³ (R² = −0.376). At +24 h, Ridge RMSE is 27.24 vs.~CAMS approximately 35.8 μg/m³ (R² = −0.184). Every locally trained model outperforms CAMS by at least 16\% in RMSE, and at 12 h the margin reaches 40\%. This consistency is of direct policy relevance: CAMS is the only freely available alternative forecast for Almaty, and the present results demonstrate that a locally trained model built on open data provides substantially superior guidance at every operational lead time. The practical recommendation is therefore unambiguous, an operational forecasting system for Almaty should use XGBoost for the 6-hour product and Ridge regression for the 12 and 24-hour products.

This reversal of model ranking with increasing horizon is the most substantive finding of the study. It is consistent with known properties of gradient-boosted tree ensembles: as the forecast horizon lengthens, the predictive signal carried by recent PM2.5 lags weakens, and the non-linear feature interactions that XGBoost exploits at short horizons become noise rather than signal. Ridge regression's L2 regularisation causes it to revert toward the mean more aggressively, which is advantageous when the target is difficult to predict and over-fitting to training-set patterns is the dominant error source. With only 9,242 training samples across 27 features, the ratio of samples to parameters is moderate, making L2 regularisation more effective than tree ensemble variance reduction at long horizons.

\subsubsection{4.1.2. RQ2, Which features drive model skill?}\label{rq2-which-features-drive-model-skill}

Hypothesis H2 posited that BLH-derived features would be among the top-five feature importances at the 6-hour horizon. The error analysis in §3.4 provides indirect evidence consistent with this hypothesis: XGBoost begins to respond to declining BLH approximately 2--3 hours before the full onset of a PM2.5 inversion event, while persistence and linear regression respond only after the observed PM2.5 has started rising. This advance response is only possible if the model has learned an association between BLH collapse and future PM2.5 increase, an association that is encoded in the BLH and blh\_delta features. Formal SHAP analysis is an identified limitation of the present work (§4.2); the claim about feature importance is therefore based on error pattern analysis rather than global attribution.

The physical interpretation of this finding is worth making explicit. If BLH is the dominant predictor of forecast skill, then the model is, in effect, learning to forecast PM2.5 from the forecast boundary-layer state: when the boundary layer is expected to compress, PM2.5 rises; when it expands, PM2.5 falls. This is a physically correct relationship, it is the mechanism by which Almaty's inversion events develop, and the model's ability to capture it is the primary source of its advantage over persistence and CAMS. CAMS, as a global chemistry-transport model operating at roughly 40 km horizontal resolution, represents the city as a single grid cell and cannot resolve the fine-scale topographic effects that drive BLH variability within the basin. A sensitivity analysis not conducted in the present study but recommended for future work would remove all BLH-derived features (BLH, blh\_delta) from the feature matrix and measure the resulting increase in RMSE. This would directly quantify the contribution of meteorological boundary-layer information to forecast skill, separating the role of BLH from that of the autoregressive PM2.5 lag features.

\subsubsection{4.1.3. RQ3, Residual error distribution and failure regimes}\label{rq3-residual-error-distribution-and-failure-regimes}

Hypothesis H3, that XGBoost under-predicts peak concentrations during inversion onset, is confirmed qualitatively by the error analysis in §3.4. During rapid inversion events (BLH collapsing from \textgreater{} 200 m to \textless{} 50 m within 3--6 hours), the XGBoost model lags the observed PM2.5 ramp by approximately 2--3 hours, which corresponds to one to two forecast time steps. At a typical inversion-onset ramp rate of 10--20 μg/m³/h, a 2-hour lag translates to an underestimation of peak concentration by 20--60 μg/m³. Given that peak concentrations during severe inversion events can reach 150--200 μg/m³, this underestimation is substantial in absolute terms but does not eliminate the model's operational value: the model detects the onset roughly 2 hours earlier than a persistence forecast, which provides a non-trivial warning lead time for issuing health advisories. A 2-hour lead time (when the model first detects the BLH-driven onset) vs.~the near-zero lead time offered by persistence remains useful for municipalities wishing to notify schools or vulnerable populations ahead of a pollution spike.

This lag structure is inherent to lag-based forecasting: the model can only respond to elevated PM2.5 once the lag features (pm25\_lag1, pm25\_lag3) encode the elevated values, which requires 1--3 hours from onset. The BLH delta feature partially compensates by providing an advance meteorological signal, but it cannot fully overcome the temporal delay imposed by the lag structure. The failure is therefore architectural rather than a consequence of under-fitting. The most direct remedy, integrating forecast rather than analysed BLH values, is described in §4.5.

Rapid clearing events reveal a symmetric failure mode: XGBoost over-predicts PM2.5 for 2--4 hours after a synoptic boundary layer breakdown because the autoregressive lags still carry the preceding high-concentration state. This over-prediction is operationally less harmful than the under-prediction failure mode, since it triggers unnecessary public health advisories rather than missing them, but it represents a quantifiable accuracy loss at the transitions that matter most for communication. The over-prediction magnitude is bounded by the same 10--20 μg/m³/h ramp rate: a 3-hour over-prediction tail produces forecast values approximately 30--60 μg/m³ above the true clearing trajectory.

\subsection{4.2. Limitations}\label{limitations}

\subsubsection{4.2.1. LCS data quality}\label{lcs-data-quality}

The PM2.5 observations underlying both the training targets and the autoregressive features are uncalibrated LCS readings. The literature on optical particle counters operating under high-humidity conditions reports hygroscopic growth biases of 20--60\% at RH \textgreater{} 80\% \textbackslash cite\{\cite{DiAntonio2018humidity, Kosmopoulos2020}\}: water droplets coating or co-accumulating with hygroscopic PM2.5 particles scatter additional light, causing the sensor to report elevated concentrations relative to the true dry-mass concentration. Almaty winter relative humidity frequently exceeds 80\% during the same inversion events that produce the highest PM2.5 concentrations. This co-occurrence is mechanistically coherent, both high humidity and high PM2.5 are products of the same stable, moist boundary-layer state, but it means that the model learns to associate high-humidity meteorological conditions with elevated sensor readings that partially reflect hygroscopic artefact rather than actual PM2.5 mass increase. The consequence is that the model's accuracy is hardest to evaluate precisely in the most operationally critical scenario: the transition into a severe inversion event.

The city-median aggregation applied across the 192-station network mitigates random inter-sensor variance but does not correct the systematic positive bias that affects all sensors simultaneously under high-humidity conditions. The absence of co-location calibration against reference instruments means that the absolute concentration values reported by the model cannot be validated against a traceable standard. The forecasting system is designed to be actionable in terms of relative concentration changes and AQI-category transitions; the reported μg/m³ values should not be treated as equivalent to reference-grade measurements. Future work extending the present pipeline should incorporate a correction factor derived from the nearest Kazhydromet reference station, as described in §4.5.

\subsubsection{4.2.2. Geographic scope and network heterogeneity}\label{geographic-scope-and-network-heterogeneity}

All results in this study apply to a single city (Almaty) over a single training window (October 2024 -- April 2026). The spatial extent of the modelling domain is a 25 km radius circle centred on the city centre (43.24°N, 76.89°E). This boundary excludes two geographically distinct emission zones that contribute to the regional PM2.5 budget. To the west of the city, a peri-urban industrial belt including cement plants, small-scale manufacturing, and solid-waste processing facilities emits particulate matter that is transported into the city under easterly wind conditions. To the east, dispersed coal-heated settlements in the Ili river delta contribute to the regional background PM2.5 during stable weather but are essentially unmonitored in the OpenAQ network. Neither emission source is represented in the training targets or features, which means the city-median target is a spatially truncated average: it captures the residential and central-district concentration field but underweights or omits industrial and peri-urban contributions. On days when easterly transport brings industrial plumes into the city, the model will lack input features encoding this source and may systematically under-predict PM2.5 even given accurate meteorological inputs.

The sensor network is spatially heterogeneous: 153 of the 192 stations are AirGradient units concentrated in central and southern residential districts; Clarity devices cover a separate set of locations; the single AirNow station is at a fixed governmental site. The city-median therefore implicitly weights central-district readings more heavily than peripheral ones during periods when peripheral stations are offline. The AirGradient firmware outage in October--November 2025 (§3.1.4), which reduced active stations from approximately 150 to 20--30, is the most prominent example: for approximately six weeks, the spatial median was effectively a central-district median, potentially biasing the training targets for that period. Predictions made during this period are less spatially representative than those in the full-coverage period.

\subsubsection{4.2.3. Absence of formal hyperparameter optimisation}\label{absence-of-formal-hyperparameter-optimisation}

The XGBoost hyperparameters used in this study were set to standard values for tabular regression problems of the present size, without a systematic grid search or Bayesian optimisation procedure. A properly optimised XGBoost model - obtained, for instance, via a Bayesian search over the space n\_estimators ∈ {[}200, 1000{]}, learning\_rate ∈ {[}0.01, 0.3{]}, and max\_depth ∈ {[}3, 10{]}, would likely achieve lower RMSE than reported here. The comparison between Ridge and XGBoost is therefore a comparison between a well-regularised linear model and a plausibly-configured but not optimised non-linear model. The reported reversal at the 12 and 24-hour horizons may shift in XGBoost's favour with improved hyperparameter selection; however, changing the sign of the reversal would require XGBoost to close gaps of 2.82 μg/m³ (at +12 h) and 4.24 μg/m³ (at +24 h) in RMSE, which would be a large gain from optimisation alone in a 9,242-row dataset. The fundamental conclusion that Ridge regression is competitive with XGBoost at longer horizons on LCS-based data is unlikely to be reversed by hyperparameter tuning.

\subsubsection{4.2.4. Temporal coverage and generalisability}\label{temporal-coverage-and-generalisability}

The training window covers parts of two winter seasons (2024--25 and 2025--26), with the 2024--25 season underrepresented because most AirGradient units were not online until late 2024. The training set contains approximately 9,242 hours across 562 days, which corresponds to approximately 1.3 heating seasons (partial 2024--25 season plus the full 2025--26 season). Published studies on gradient boosting for urban air quality suggest that 2--3 full annual cycles are needed for stable generalisation \cite{enebish2021ulaanbaatar, menares2021forecasting}: a single year of training captures one realisation of the synoptic-scale variability that drives inversion frequency and severity, while two or more years provide the model with a range of winter conditions against which to generalise. The present model has been trained on 1.3 cycles. Year-to-year meteorological variability, particularly the frequency and severity of temperature-inversion episodes, may mean that the reported test-set performance overestimates or underestimates the model's accuracy in a future winter season with different synoptic patterns. As the AirGradient network accumulates additional years of data, periodic retraining is expected to improve generalisation and should be built into any operational deployment protocol.

The summer months are also underrepresented in the training data relative to their share of a full annual cycle. The training window of October 2024 -- April 2026 captures two winters but only one summer (June--September 2025), during which PM2.5 monthly means ranged from approximately 12 to 17 μg/m³ - well below the WHO guideline and largely driven by dust rather than combustion. A model trained predominantly on winter data may therefore overweight the inversion-onset signals and underweight the dust transport features that govern summer-season PM2.5 exceedances. This seasonal imbalance is an additional argument for extending the training window to cover at least two complete annual cycles before treating the model as suitable for year-round operational deployment.

\subsection{4.3. Comparison with related approaches}\label{comparison-with-related-approaches}

The results reported here are broadly consistent with the published literature on gradient-boosting PM2.5 forecasters in comparable settings. Studies comparing XGBoost to persistence, ARIMA, and physics-based models in Asian urban environments with strong inversion seasonality typically report R² values of 0.45--0.65 at 6-hour horizons and 0.30--0.50 at 24-hour horizons for city-level daily or hourly targets \cite{enebish2021ulaanbaatar, qi2019pm25lstm}. The present study's 6-hour R² = 0.545 falls within the central range of these published values. The 24-hour R² = 0.175 from XGBoost is below the published lower bound of 0.30; however, Ridge regression achieves R² = 0.383 at the same horizon, which falls within the expected range. The gap between the XGBoost and Ridge 24-hour values (0.175 vs.~0.383) suggests that the shortfall relative to the literature for XGBoost reflects over-fitting on the moderate-size training set rather than a fundamental inadequacy of the feature set. The use of uncalibrated LCS readings as the training target introduces additional noise in the target variable that would reduce R² at all horizons relative to a model trained on reference-grade data, so the comparison to reference-grade-trained published benchmarks should be read with this caveat.

The CAMS comparison (R² ≤ −0.184 at all horizons, with values of −0.314 at +6 h, −0.376 at +12 h, and −0.184 at +24 h) is consistent with expectations for a roughly 40 km global CTM applied to an inversion-dominated mountain basin. At that grid resolution, the model does not resolve the bowl topography that governs BLH suppression in Almaty, and the single grid cell covering the city cannot distinguish the inversion regime from the surrounding steppe. Published studies applying CAMS output to Central Asian cities with inversion regimes comparable to Almaty's report similarly poor performance, motivating the use of local statistical post-processing as a required correction layer \cite{inness2019cams, Wang2016CTM}. The present work demonstrates that even a simple Ridge regression model trained on LCS data substantially outperforms CAMS without requiring a separate bias-correction post-processing step.

The broader post-Soviet and Central Asian urban context warrants direct comment. Bishkek (Kyrgyzstan) operates under a comparable inversion regime, with coal-fired district heating and a similar mountain-basin topography; as of the writing of this thesis, no peer-reviewed ML-based PM2.5 forecast model calibrated to Bishkek's sensor network appears in the published literature \textbackslash cite\{\cite{Isaev2022}\}. Ulaanbaatar (Mongolia) is the most studied inversion-dominated Central/Inner Asian capital, with several published AQ modelling studies using MODIS aerosol optical depth and ground-based observations \textbackslash cite\{\cite{enebish2021ulaanbaatar}\}, though open reproducible LCS-trained forecast models remain absent from the published record. Tashkent (Uzbekistan) has published air quality trend analyses but no gradient-boosting forecast models in the peer-reviewed literature. The absence of published baselines for these cities means that the present work is positioned as an initiating study in the region rather than an incremental improvement over existing work.

One feature of the present study that distinguishes it from most published comparisons is the explicit inclusion of BLH as a predictor. The majority of published gradient-boosting PM2.5 forecasters for Asian cities use surface meteorological variables (temperature, wind speed, relative humidity, atmospheric pressure) but not BLH, which requires either a radiosonde programme or a reanalysis product. The Open-Meteo ERA5-derived BLH archive \cite{zippenfenig_openmeteo_2023} provides this variable at no cost and without registration, which is the enabling condition for including it here. The advance response to BLH collapse observed in §4.1.2 suggests that this variable carries predictive information not available in surface fields alone, a finding that would be worth replicating in future work on other mountain-basin cities where the same reanalysis product is accessible.

\subsection{4.4. Practical and scientific contributions}\label{practical-and-scientific-contributions}

The practical contribution of this work is a complete, open-source PM2.5 forecasting pipeline for Almaty that any researcher or municipal agency can deploy on free-tier infrastructure. The pipeline requires only a personal OpenAQ API key (freely available) and runs on a standard laptop or a Streamlit Community Cloud instance without GPU or cloud database dependencies. The three trained models (6, 12, and 24-hour horizon) are serialised files that can be updated by re-running the evaluation script as new sensor data accumulates. The Streamlit application provides a map-based visualisation of real-time sensor readings alongside the 6, 12, and 24-hour forecasts, accessible without registration from any device with a browser. The 30--60 second refresh cycle and absence of authentication make it directly usable by schools, hospitals, and municipal public health offices without technical intermediaries.

The first scientific contribution is the first published evaluation of XGBoost vs.~CAMS for PM2.5 forecasting in Almaty. The quantitative result, that locally trained models outperform CAMS by 16--40\% in RMSE across all three horizons, establishes the value of local model training over global CTM application in this specific geographic and meteorological context. This finding is directly relevant to the question of whether Kazhydromet or other regional agencies should invest in local model development as a complement or replacement to CAMS-based guidance, and the answer supported by the present data is clearly affirmative.

The second scientific contribution is the horizon-dependent model-complexity finding. The result that Ridge regression outperforms XGBoost at the 12 and 24-hour horizons on a training set of 9,242 samples with 27 features adds a data point to the ongoing methodological discussion about model selection for air quality forecasting with LCS data. The finding is consistent with the theoretical argument that L2 regularisation is more effective than ensemble variance reduction when the target is noisy and the sample count is moderate, but this has not been demonstrated previously on an LCS-based Central Asian dataset. This result has direct design implications: operational forecasting systems in similar data environments should not default to gradient-boosted ensembles without testing regularised linear alternatives at longer horizons.

A third contribution, less often recognised in ML-focused theses, is the documentation of the sensor-network operational behaviour over a 562-day window. The firmware outage of October--November 2025 that reduced the active station count from approximately 150 to 20--30 for roughly six weeks is recorded in the data pipeline and its effect on the spatial representativeness of the target variable is discussed in §4.2.2. Documenting outage events, their duration, and their downstream effect on training-data quality is a necessary step for any group seeking to build on this dataset in future work. The SQLite database and the associated data-quality flags in \texttt{data/almaty\_aq.db} constitute a secondary research output that has value independent of the forecasting models themselves, providing a curated 562-day record of Almaty PM2.5 at hourly resolution that is not available from any other open source in this form.

\subsection{4.5. Future directions}\label{future-directions}

Five extensions are identified as high-priority for future work.

\textbf{SHAP-based feature attribution.} A formal SHAP analysis, computing Shapley values for each feature across the full test set, would directly answer RQ2 and provide a quantitative, peer-reviewable attribution of model skill to BLH, autoregressive lags, and other feature families. The \texttt{shap.TreeExplainer} class computes exact Shapley values for XGBoost in O(TLD²) time, where T is the number of trees, L is the number of leaves per tree, and D is tree depth; for the present models (approximately 300 trees, depth 6, and 27 features), this computation completes in seconds on a laptop. The analysis was deferred from the present study due to time constraints and will be included in the published version of this work.

\textbf{Forecast BLH integration.} The Open-Meteo forecast API (\texttt{api.open-meteo.com/v1/forecast}, parameters \texttt{hourly=boundary\_layer\_height}) delivers 10-day-ahead BLH predictions from the ECMWF IFS model at the same spatial resolution as the historical archive used in this study. Replacing the historical analysed BLH feature with the corresponding forecast BLH value at each prediction step would allow the model to look ahead on the boundary-layer trajectory rather than relying solely on the current analysed state. This single change would be expected to reduce the 2--3 hour onset lag identified in §4.1.3, since the model would receive advance notice of an impending BLH collapse approximately 12--24 hours before it occurs. Integration requires only a change to the feature assembly step in \texttt{src/features/build.py} to query the forecast endpoint rather than the archive endpoint.

\textbf{LCS bias correction via reference co-location.} The most impactful data-quality improvement would be a short (2--4 week) co-location campaign deploying a reference-grade BAM-1020 or TEOM instrument adjacent to one of the AirGradient stations in central Almaty, combined with a humidity-dependent correction factor (such as the Kano correction or a quadratic RH regression) estimated from the paired data. A four-week co-location spanning a mix of high- and low-humidity conditions would provide enough paired observations to estimate correction coefficients that could then be applied across the full network. This would make the absolute concentration estimates more reliable and is expected to reduce the systematic over-prediction bias at high-humidity inversion events.

\textbf{Extension to Shymkent and Karaganda.} Both cities share Almaty's LCS-on-OpenAQ data profile and seasonally driven PM2.5 regime. The pipeline architecture is fully parameterised by city centre coordinates and API window: changing the \texttt{CITY\_LAT}, \texttt{CITY\_LON}, and \texttt{CITY\_RADIUS\_KM} parameters in \texttt{src/config.py}, then re-running the data ingest and model training scripts, constitutes the complete adaptation procedure. No code changes beyond the configuration file are required. This makes geographic extension tractable for researchers or municipalities in other Central Asian cities without requiring specialised software development.

\textbf{Ensemble of XGBoost and Ridge regression.} Since XGBoost produces the lowest RMSE at the 6-hour horizon and Ridge regression produces the lowest RMSE at the 12 and 24-hour horizons, a blended ensemble using horizon-dependent weights is a natural next step. A simple weighted average - with weights learned by minimising RMSE on a held-out validation set for each horizon separately, would be expected to match or exceed both individual models at all three horizons. At the 6-hour horizon the optimal ensemble would weight XGBoost heavily; at 24 hours it would weight Ridge heavily; and at 12 hours the appropriate blend is an empirical question. This ensemble requires no new feature engineering or model architecture work, only the addition of a weight optimisation step in the evaluation pipeline, and is a low-cost path to marginal accuracy improvements at all operational horizons.

Beyond these five targeted extensions, a longer-term research direction is the incorporation of traffic emission proxies into the feature matrix. Almaty's PM2.5 is dominated by coal combustion during winter inversions, but vehicular emissions contribute meaningfully to the annual-mean concentration, particularly during morning and evening rush hours on days when the boundary layer is partially but not fully suppressed (BLH in the range 50--200 m). Hourly vehicle counts are not available from open sources for Almaty, but proxy variables - such as hour-of-day interaction terms with working-day indicators, could partially capture the traffic signal. Adding four to six such interaction features to the existing 27-feature matrix would require no change to the model architecture and could reduce residuals in the moderate-concentration regime (20--60 μg/m³) where the current models show their largest proportional errors.
